\chapter*{PHIẾU ĐÁNH GIÁ VÀ NHẬN XÉT CỦA GIẢNG VIÊN HƯỚNG DẪN}
\addcontentsline{toc}{chapter}{PHIẾU ĐÁNH GIÁ CỦA GIẢNG VIÊN HƯỚNG DẪN}

\noindent \textbf{Trường:} TRƯỜNG ĐẠI HỌC CÔNG NGHỆ THÔNG TIN VÀ TRUYỀN THÔNG \\
\textbf{Học phần:} Ứng dụng trí tuệ nhân tạo (AIA331) \quad | \quad \textbf{Mã số:} 80273 \\
\textbf{Đề tài:} HỆ THỐNG QUẢN LÝ KHO CÓ TÍCH HỢP AI \\
\textbf{Giảng viên hướng dẫn:} ThS. Nguyễn Tuấn Anh \\
\textbf{Sinh viên thực hiện:} 1. Dương Minh Quang (DTC245200225) \quad 2. Vũ Ngọc Anh Thuyên (DTC245200428)

\vspace{4mm}
\begin{longtable}{|c|p{8.5cm}|c|c|}
\hline
\rowcolor{tablehead}
\textbf{STT} & \textbf{Nội dung tiêu chí đánh giá (Rubric AIA331)} & \textbf{Vị trí} & \textbf{Đánh giá} \\ \hline
\endhead
01 & Khảo sát bối cảnh, đối tượng người dùng, quy trình nghiệp vụ kho & Chương 1, 3 & Đạt \\ \hline
02 & Đặc tả các yêu cầu chức năng (REQ-AUTH -> REQ-AUDIT) \& phi chức năng & Chương 3 & Đạt \\ \hline
03 & Phân tích quy trình hệ thống qua Use Case Diagram chuẩn UML & Chương 3, Hình 3.1 & Đạt \\ \hline
04 & Mô hình hóa nghiệp vụ Đăng nhập qua Activity Diagram & Chương 4, Hình 4.1 & Đạt \\ \hline
05 & Mô hình hóa nghiệp vụ Nhập kho qua Activity Diagram & Chương 4, Hình 4.2 & Đạt \\ \hline
06 & Mô hình hóa nghiệp vụ Xuất kho chống âm qua Activity Diagram & Chương 4, Hình 4.3 & Đạt \\ \hline
07 & Mô hình hóa luồng xử lý AI qua Activity Diagram & Chương 4, Hình 4.4 & Đạt \\ \hline
08 & Thiết kế Cơ sở dữ liệu quan hệ chuẩn 3NF (11 bảng) & Chương 5, Bảng 5.1-5.11 & Đạt \\ \hline
09 & Sơ đồ thực thể liên kết ERD chuẩn hóa & Chương 5, Hình 5.1 & Đạt \\ \hline
10 & Khóa chính, khóa ngoại, ràng buộc toàn vẹn và Indexing tối ưu & Chương 5, ddl.sql & Đạt \\ \hline
11 & Mô hình kiến trúc tổng thể 3-Tier Client-Server & Chương 6, Hình 6.1 & Đạt \\ \hline
12 & Sơ đồ luồng dữ liệu DFD Mức 0 (Context Diagram) & Chương 6, Hình 6.2a & Đạt \\ \hline
13 & Sơ đồ luồng dữ liệu DFD Mức 1 (Chi tiết phân rã nghiệp vụ) & Chương 6, Hình 6.2b & Đạt \\ \hline
14 & Sơ đồ tuần tự Sequence Diagram cho quy trình Nhập kho & Chương 7, Hình 7.1 & Đạt \\ \hline
15 & Sơ đồ tuần tự Sequence Diagram cho quy trình Xuất kho & Chương 7, Hình 7.2 & Đạt \\ \hline
16 & Sơ đồ tuần tự Sequence Diagram cho quy trình Hỏi đáp AI NLP & Chương 8, Hình 8.1 & Đạt \\ \hline
17 & Thuật toán tính toán tồn an toàn và tốc độ tiêu thụ (Burn-rate) & Chương 8, Công thức (8.1) & Đạt \\ \hline
18 & Xử lý an toàn ngoại lệ chia cho 0 (ZeroDivisionError) khi sales=0 & Chương 8, Mã nguồn 8.1 & Đạt \\ \hline
19 & Cơ chế Strict Grounding Context chống ảo giác cho mô hình ngôn ngữ & Chương 8, Bảng 8.3 & Đạt \\ \hline
20 & Cơ chế Fallback tự động sang Rule Engine khi mất kết nối LLM & Chương 8, Bảng 8.4 & Đạt \\ \hline
21 & Kiến trúc Backend FastAPI, Pydantic Schemas và SQLAlchemy ORM & Chương 9 & Đạt \\ \hline
22 & Thiết kế hệ thống RESTful API chuẩn HTTP status codes & Chương 9, Bảng 9.1 & Đạt \\ \hline
23 & Quản lý phiên và phân quyền RBAC (ADMIN, WAREHOUSE, ACCOUNTANT) & Chương 9, 11 & Đạt \\ \hline
24 & Giao diện người dùng React Vite Responsive hiện đại & Chương 10, Hình 10.1-10.5 & Đạt \\ \hline
25 & Bảng điều khiển KPI Dashboard trực quan hóa số liệu thời gian thực & Chương 10, Hình 10.2 & Đạt \\ \hline
26 & Giao diện AI Hub phân tích báo cáo kho tự động & Chương 10, Hình 10.3 & Đạt \\ \hline
27 & Giao diện Trợ lý thông minh hỏi đáp bằng tiếng Việt tự nhiên & Chương 10, Hình 10.4 & Đạt \\ \hline
28 & Cơ chế bảo mật JWT Bearer Token với thuật toán mã hóa HS256 & Chương 11 & Đạt \\ \hline
29 & Mã hóa mật khẩu một chiều sử dụng thuật toán Argon2/Bcrypt & Chương 11 & Đạt \\ \hline
30 & Giao dịch CSDL đảm bảo tính ACID và Rollback khi xuất vượt tồn & Chương 11, Bảng 11.2 & Đạt \\ \hline
31 & Ghi nhật ký kiểm toán (Audit Trail) cho mọi biến động xuất nhập tồn & Chương 11 & Đạt \\ \hline
32 & Cấu hình biến môi trường an toàn qua .env / .env.example & Chương 12 & Đạt \\ \hline
33 & Quy trình đóng gói Docker container và triển khai hệ thống & Chương 12 & Đạt \\ \hline
34 & Chiến lược kiểm thử tự động Unit Test và Integration Test & Chương 13 & Đạt \\ \hline
35 & Bộ kịch bản kiểm thử 12 Test Cases bao phủ chức năng chính & Chương 13, Bảng 13.1 & Đạt \\ \hline
36 & Kết quả thực thi Pytest 100% Passed (12/12 test cases) & Chương 14, Bảng 14.1 & Đạt \\ \hline
37 & Đo lường định lượng độ trễ API (<100ms) và phản hồi AI (<1.5s) & Chương 14, Bảng 14.2 & Đạt \\ \hline
38 & Phân loại minh bạch 5 hạng mục: IMPLEMENTED, TESTED, PROPOSED... & Chương 15, Bảng 15.1 & Đạt \\ \hline
39 & Ma trận truy vết yêu cầu (Traceability Matrix) từ Req đến Test & Phụ lục A, Bảng A.1 & Đạt \\ \hline
40 & Định dạng báo cáo chuẩn ICTU (Font Times, lề, khung TikZ, caption) & Toàn bộ báo cáo & Đạt \\ \hline
\end{longtable}

\vspace{4mm}
\noindent \textbf{Nhận xét chung của Giảng viên hướng dẫn:} \\
\noindent \rule{\textwidth}{0.4pt} \\
\noindent \rule{\textwidth}{0.4pt}

\vspace{3mm}
\noindent \textbf{Điểm số:} .............. / 10.0 \hfill \textbf{Điểm chữ:} ....................................

\vspace{8mm}
\hfill
\begin{tabular}{c}
Thái Nguyên, ngày ..... tháng ..... năm 2026 \\
\textbf{GIẢNG VIÊN HƯỚNG DẪN} \\[2.5cm]
\textbf{ThS. NGUYỄN TUẤN ANH}
\end{tabular}
